% ==============================================================
% Der Explanatory Gap -- Bestandsaufnahme, Loesungsversuche
% und die Position des Ontophaenomenalismus
% Dokument EG-1, Version 1.9
% Kompiliert mit: pdfLaTeX
% ==============================================================

\documentclass[11pt,a4paper]{article}

\usepackage[utf8]{inputenc}
\usepackage[T1]{fontenc}
\usepackage[german]{babel}
\usepackage{amsmath,amssymb,amsthm}
\usepackage{geometry}
\usepackage{parskip}
\usepackage{enumitem}
\usepackage{booktabs}
\usepackage{xcolor}
\usepackage{hyperref}

\geometry{margin=2.5cm}
\setlist{itemsep=0.2em, topsep=0.2em}
\hypersetup{
    colorlinks=true,
    linkcolor=blue,
    urlcolor=blue,
    pdftitle={Der Explanatory Gap -- Dokument EG-1, Version 1.9},
    pdfauthor={Forschungsprogramm Ontoph\"{a}nomenalismus}
}

\theoremstyle{plain}
\newtheorem{proposition}{Proposition}[section]
\newtheorem{corollary}{Korollar}[section]
\newtheorem{theorem}{Theorem}[section]

\theoremstyle{definition}
\newtheorem{definition}{Definition}[section]
\newtheorem{remark}{Bemerkung}[section]
\newtheorem{methodennote}{Methodennote}[section]

\title{
  \textbf{Der Explanatory Gap}\\[0.4em]
  \large Bestandsaufnahme, L\"{o}sungsversuche und die Position\\
  des Ontoph\"{a}nomenalismus\\[0.5em]
  {\normalsize Dokument EG-1, Version 1.9}
}
\author{
  \textbf{Forschungsprogramm Ontoph\"{a}nomenalismus}\\
  \small Siehe TRANSPARENCY f\"{u}r methodologische Urheberschaft und KI-Einsatz.
}
\date{Juni 2026}

\begin{document}

\maketitle

\begin{abstract}
Der \emph{Explanatory Gap} --- die L\"{u}cke zwischen physikalischer oder struktureller
Beschreibung und ph\"{a}nomenalem Erleben --- ist das tiefste und zugleich universellste
Problem der Bewusstseinsphilosophie. Kein bestehendes System hat es geschlossen. Der
Ontoph\"{a}nomenalismus (OP) beansprucht nicht, es zu l\"{o}sen. Er beansprucht etwas
Bescheideneres und Ehrlicheres: den Gap pr\"{a}zise zu \emph{relokieren} und zu zeigen,
dass er in seiner klassischen Formulierung m\"{o}glicherweise auf einer \emph{kategorialen
Fehlfrage} beruht.

Es gelten die methodologischen Konventionen gem\"{a}\ss{} MKO. Das Fragedifferential
wird vollst\"{a}ndig im Dokument DFD (Das Fragedifferential) behandelt; EG-1 verweist
darauf und verwendet die wesentlichen Ergebnisse.

Dieses Dokument kartografiert in acht Schritten: erstens, was der Gap genau ist;
zweitens, welche Ans\"{a}tze ihn zu schlie\ss{}en versucht haben und wo sie scheitern;
drittens, was alle Ans\"{a}tze strukturell gemeinsam haben; viertens, die Symmetrisierung
des Hard Problems --- eine zweistufige Dekonstruktion der materialistischen Asymmetrie;
f\"{u}nftens, ob die klassische Formulierung des Gaps selbst eine unbegr\"{u}ndete
Pr\"{a}misse enth\"{a}lt; sechstens, welche spezifische Stellung der OP einnimmt;
siebtens, das Fragedifferential als Individuationsprinzip; achtens, die
erkenntnistheoretischen Konsequenzen.
\end{abstract}

\tableofcontents
\newpage

%==============================================================
\section{Was der Explanatory Gap ist --- und was nicht}
%==============================================================

\subsection{Die Formulierung des Problems}

Der Begriff \emph{Explanatory Gap} wurde 1983 von Joseph Levine eingef\"{u}hrt, um ein
Ph\"{a}nomen zu beschreiben, das Philosophen und Wissenschaftler seit Descartes beunruhigt:
Selbst wenn wir eine vollst\"{a}ndige physikalische oder funktionale Beschreibung eines
Gehirnzustands besitzen, bleibt eine Frage offen, die durch diese Beschreibung nicht
beantwortet wird:

\begin{quote}
\emph{Warum f\"{u}hlt es sich so an?}
\end{quote}

Warum ist das Feuern bestimmter Neuronen nicht einfach ein Informationsverarbeitungsprozess
--- warum ist es \emph{schmerzhaft}? Warum ist die Wellenl\"{a}nge 700\,nm nicht einfach
ein physikalischer Parameter --- warum sieht man \emph{Rot}? Die physikalische Beschreibung
mag vollst\"{a}ndig und korrekt sein; und dennoch scheint etwas zu fehlen: die Erkl\"{a}rung,
warum dieser Prozess \emph{\"{u}berhaupt etwas ist, wie es ist}.

David Chalmers hat dieses Problem 1995 als das \emph{Hard Problem of Consciousness}
reformuliert und von den \glqq{}Easy Problems\grqq{} abgegrenzt --- jenen Fragen \"{u}ber
kognitive Funktionen, Verhalten und Informationsverarbeitung, die prinzipiell durch
Neurowissenschaft und kognitive Psychologie beantwortbar sind. Das Hard Problem ist das
einzige, das auch nach vollst\"{a}ndiger funktionaler Erkl\"{a}rung offen bleibt.

\begin{definition}[Explanatory Gap --- Arbeitsdefinition]
Der Explanatory Gap bezeichnet die strukturelle L\"{u}cke zwischen jeder Beschreibung
eines Systems in objektiven, drittpersonalen Begriffen (physikalisch, funktional,
informationstheoretisch, mathematisch) und der Tatsache, dass dieses System subjektives
Erleben \emph{hat} --- dass es \glqq{}etwas ist, wie es ist\grqq{} (Nagel), dieses System
zu sein. Die L\"{u}cke ist nicht eine des Wissens, sondern eine der \emph{Erkl\"{a}rung}:
Kein Zuwachs an objektiven Informationen scheint sie prinzipiell schlie\ss{}en zu
k\"{o}nnen.
\end{definition}

\subsection{Warum der Gap anders ist als andere Erkl\"{a}rungsprobleme}

\textbf{Reduktive Erkl\"{a}rungen funktionieren} in der Wissenschaft \"{u}berall: Temperatur
wird erkl\"{a}rt als mittlere kinetische Energie von Molek\"{u}len; Leben als bestimmte
Art chemischer Selbstorganisation; Gene als Informationssequenzen in DNA. In allen diesen
F\"{a}llen gilt: Wenn wir die strukturelle Erkl\"{a}rung haben, \emph{verschwindet} das
Erkl\"{a}rungsproblem. Es gibt kein Residuum.

\textbf{Beim Bewusstsein bleibt ein Residuum.} Selbst eine vollst\"{a}ndige
neurowissenschaftliche Erkl\"{a}rung des Schmerzes --- Nozizeptoren, C-Fasern,
Thalamuskortex-Schleifen, Aussch\"{u}ttung von Substanz~P --- l\"{a}sst offen, warum
dieses Netz von Prozessen \emph{wehtut}. Ein philosophischer Zombie (Chalmers) --- ein
System, das all diese Prozesse vollst\"{a}ndig nachahmt, ohne jedes Erleben --- ist
konzeptuell koh\"{a}rent. Das bedeutet: Erleben folgt logisch nicht aus Funktion.

\begin{methodennote}[Die Asymmetrie des Gaps]
Der Explanatory Gap ist asymmetrisch: Es gibt keine symmetrische L\"{u}cke in die andere
Richtung. Wenn ich Schmerz erlebe, zweifle ich nicht daran, dass ein physikalischer
Prozess stattfindet. Die L\"{u}cke geht nur in eine Richtung --- von der objektiven
Beschreibung zum subjektiven Erleben, nicht umgekehrt. Diese Asymmetrie ist selbst
erkl\"{a}rungsbed\"{u}rftig und bildet einen Ausgangspunkt f\"{u}r den OP.
\end{methodennote}

\subsection{Qualia und das \glqq{}Was-es-ist\grqq{}}

Thomas Nagel hat 1974 in seinem Aufsatz \glqq{}What Is It Like to Be a Bat?\grqq{} die
subjektive Dimension des Bewusstseins in eine pr\"{a}gnante Formel gefasst: Ein System
ist bewusst genau dann, wenn es \emph{etwas ist, wie es ist}, dieses System zu sein.
Diese subjektiven Erlebensqualit\"{a}ten werden als \emph{Qualia} bezeichnet: Das Rot von
Rot, das Wehtun von Schmerz, das Klingen von Musik. Sie sind das, was nach jeder
funktionalen Erkl\"{a}rung \"{u}brigbleibt.

%==============================================================
\section{L\"{o}sungsversuche und ihre Scheiterstellen}
%==============================================================

\subsection{Eliminativer Materialismus}
\textbf{Position:} Es gibt keine Qualia (Churchland, Dennett).
\textbf{Scheiterstelle:} Performativ inkoh\"{a}rent --- die Leugnung selbst ist ein
Erleben.

\subsection{Identit\"{a}tstheorie und Typ-Physikalismus}
\textbf{Position:} Mentale Zust\"{a}nde \emph{sind} neuronale Zust\"{a}nde (Smart,
Armstrong).
\textbf{Scheiterstelle:} Setzt voraus, was erkl\"{a}rt werden soll. Warum ist
C-Faser-Feuern schmerzhaft? Unbeantwortet. Kripkes Modalargument versch\"{a}rft dies.

\subsection{Funktionalismus}
\textbf{Position:} Mentale Zust\"{a}nde werden durch kausale Rollen definiert (Putnam).
\textbf{Scheiterstelle:} Das Zombie-Argument: Funktionale Rollen k\"{o}nnen ohne Erleben
realisiert werden. Ned Block unterscheidet \emph{Access Consciousness} von
\emph{Phenomenal Consciousness}: Der Funktionalismus erkl\"{a}rt ersteres, nicht letzteres.

\subsection{Eigenschaftsdualismus und Russellianischer Monismus}
\textbf{Position:} Ph\"{a}nomenale Eigenschaften sind irreduzibel (Chalmers); intrinsische
Eigenschaften physikalischer Dinge sind proto-ph\"{a}nomenal (Maxwell, Chalmers).
\textbf{Scheiterstelle:} Beschreibt den Gap, ohne ihn zu schlie\ss{}en. Das Kombinations-
problem bleibt ungelöst.

\subsection{Panpsychismus und Kosmopsychismus}
\textbf{Position:} Erleben ist fundamental (Goff, Strawson); das Universum hat Erleben
(Kosmopsychismus).
\textbf{Scheiterstelle:} Kombinationsproblem (Panpsychismus) bzw.\ Differenzierungs-
problem (Kosmopsychismus) --- letzteres ist genau das zentrale Desiderat des OP.

\subsection{Integrierte Informationstheorie (IIT)}
\textbf{Position:} Bewusstsein ist identisch mit $\Phi$ (Tononi).
\textbf{Scheiterstelle:} Identifiziert, erkl\"{a}rt nicht. Einfache Gitter-Netzwerke
haben h\"{o}heres $\Phi$ als das menschliche Gehirn (Aaronson) --- kontra-intuitiv.

\subsection{Prozessphilosophie: Whitehead}
\textbf{Position:} Jedes \glqq{}actual occasion\grqq{} hat eine erlebende Seite
(\emph{prehension}).
\textbf{Scheiterstelle:} L\"{o}st den Gap durch metaphysische Ausweitung --- keine
Erkl\"{a}rung. Die qualitative Differenz zwischen Elektron-Pr\"{a}hension und
menschlichem Schmerz bleibt ungekl\"{a}rt.

\subsection{Metzinger: Selbstmodell und Transparenz}
\textbf{Position:} Ph\"{a}nomenales Bewusstsein entsteht durch ein transparentes
Selbstmodell.
\textbf{Scheiterstelle:} Erkl\"{a}rt nicht, warum Transparenz \emph{Erleben} erzeugt.
Der Gap tritt neu auf: Warum f\"{u}hlt sich Transparenz nach etwas an?

%==============================================================
\section{Was alle Ans\"{a}tze gemeinsam haben}
%==============================================================

\begin{proposition}[Universalit\"{a}t des Gaps]
Kein bisher formulierter Ansatz schlie\ss{}t den Explanatory Gap. Alle Strategien fallen
in eine von drei Kategorien:
\begin{enumerate}[label=(\alph*)]
  \item \textbf{Leugnung:} Performativ inkoh\"{a}rent (Eliminativismus).
  \item \textbf{Umbenennung:} Neuer Begriff, keine Erkl\"{a}rung (Identit\"{a}tstheorie,
        IIT).
  \item \textbf{Verschiebung:} Problem tritt ver\"{a}ndert wieder auf (Panpsychismus,
        Whitehead, Russellianischer Monismus).
\end{enumerate}
\end{proposition}

\begin{remark}[Der Gap als strukturelles Merkmal]
Der Explanatory Gap ist kein Defekt eines bestimmten Ansatzes, sondern ein strukturelles
Merkmal des Problems selbst. Er entsteht, weil wir eine Erste-Person-Tatsache (Erleben)
durch Dritte-Person-Begriffe (Struktur, Funktion, Information) erkl\"{a}ren wollen. Die
entscheidende Frage ist: \glqq{}Welcher Ansatz ist am ehrlichsten \"{u}ber seine
Unf\"{a}higkeit, ihn zu schlie\ss{}en?\grqq{}
\end{remark}

%==============================================================
\section{Die Symmetrisierung des Hard Problems}
%==============================================================

Bevor wir die Position des OP entfalten, ist eine strategische Vorbemerkung erforderlich.
Das Hard Problem leidet unter einer tiefsitzenden methodologischen Asymmetrie --- einer
unbewussten Vorentscheidung zugunsten der materialistischen Ausgangsposition. Der OP
begegnet diesem Dogma mit einer zweistufigen Dekonstruktion.

\subsection{Stufe 1: Die empirische Symmetrie und die unbegr\"{u}ndete Zombie-Hypothese}

Die materialistische Standardfrage lautet: \emph{\glqq{}Wie und warum entsteht aus einer
hinreichend komplexen physikalischen Struktur ($A$) subjektives Erleben ($B$)?\grqq{}}
Diese Formulierung setzt stillschweigend voraus, dass eine physikalische Struktur ohne
ph\"{a}nomenale Innenseite existieren k\"{o}nnte --- die Hypothese des \emph{philosophischen
Zombies}. Diese Annahme ist jedoch selbst eine starke metaphysische Hypothese, keine
neutrale Ausgangslage.

Unsere gesamte empirische Datenbasis zeigt ausnahmslos einen bikonditionales Befund:
\[
  A \iff B
\]
wobei $A$ f\"{u}r die hochkomplexe physikalische Organisation steht und $B$ f\"{u}r
ph\"{a}nomenales Erleben. Es gibt in der Wissenschaftsgeschichte keinen einzigen
dokumentierten Fall, in dem $A$ ohne $B$ oder $B$ ohne $A$ isoliert aufgetreten w\"{a}re.

\begin{remark}[Symmetrisierung der Beweislast]
Die Beweislast ist exakt symmetrisch: Wer fordert, das Auftreten von $B$ bei $A$ m\"{u}sse
erkl\"{a}rt werden \emph{weil} $A$ ja auch ohne $B$ m\"{o}glich sein k\"{o}nnte, stellt
eine unbegr\"{u}ndete asymmetrische Forderung. Es ist logisch ebenso valide anzunehmen,
dass $A$ ohne $B$ unm\"{o}glich ist --- \emph{analog zu} (MKO, Ebene~4) der Tatsache,
dass Feuer nicht ohne W\"{a}rme existieren kann. Beide Positionen stehen auf exakt
derselben empirischen Basis.

Der Materialist muss daher genauso begr\"{u}nden, warum $A$ \"{u}berhaupt ohne $B$
m\"{o}glich sein sollte, wie der Ph\"{a}nomenologe das Auftreten von $B$ erkl\"{a}ren
muss. Das Hard Problem ist in seiner klassischen Formulierung kein neutraler Ausgangspunkt,
sondern bereits eine materialistisch pr\"{a}judizierte Fragestellung.
\end{remark}

\subsection{Stufe 2: Projektion statt einseitiger Emergenz --- das Kugel-Schatten-Prinzip}

Wenn die Asymmetrie der Ausgangsposition aufgebrochen ist, erlaubt die gewonnene
Symmetrie eine fundamentale Umkehrung der postulierten Erkl\"{a}rungsrichtung. Statt der
materialistischen Richtung ($A \to B$) betrachten wir die ph\"{a}nomenologische Richtung:
\[
  B \to A
\]

Das unhintergehbare, prim\"{a}r Gegebene ist das ph\"{a}nomenale Erleben ($B$). Gem\"{a}\ss{}
dem Denk-Monismus des OP ist die ph\"{a}nomenale Realit\"{a}t der Urgrund. Aus der
Existenz dieses Urgrundes folgt seine physikalische Manifestation ($A$) --- nicht als
kausale Emergenz, sondern als projektive Erscheinungsform derselben ontischen Struktur.

\begin{remark}[Das Kugel-Schatten-Prinzip --- Analogie, MKO Ebene~4]
Eine Kugel wirft einen Schatten an die Wand. Der Schatten ist eine reale, mathematisch
beschreibbare Repr\"{a}sentation der Kugel --- jedoch in einer dimensional reduzierten
Projektion. Die Kugel erzeugt ihren Schatten nicht im gew\"{o}hnlichen kausalem Sinne,
ebenso wenig erzeugt der Schatten die Kugel. Beide sind unterschiedliche Projektionen
derselben zugrundeliegenden Struktur unter verschiedenen Projektionsbedingungen.

Analog k\"{o}nnte gelten: Das physische Gehirn ($A$) ist nicht der Generator von
Bewusstsein, sondern die \emph{zwingende physikalische Manifestation} --- der Schatten
--- eines aktiven, strukturierten Denkflusses ($B$). Wer fragt, wie das Gehirn Bewusstsein
erzeugt, begeht einen Kategorienfehler: Er fragt, wie der Schatten die Kugel erzeugt.

Die Projektive Reduktion (DPR) w\"{u}rde in diesem Verst\"{a}ndnis nicht erkl\"{a}ren,
wie aus Physik Erleben entsteht oder umgekehrt, sondern wie eine gemeinsame ontische
Struktur (OPP, Denkfeld) unter unterschiedlichen Projektionsbedingungen einmal als
ph\"{a}nomenale Realit\"{a}t und einmal als physikalische Beschreibung erscheint.
\end{remark}

\begin{methodennote}[Strategischer Status dieser Sektion --- MKO Ebene~2]
Die Symmetrisierung des Hard Problems ist kein neues Dogma. Der OP behauptet nicht
absolut: \glqq{}Es ist so und nicht anders.\grqq{} Er formuliert die Umkehrung
($B \to A$) als ein logisch vollkommen gleichwertiges M\"{o}glichkeitsszenario. Dem
Materialismus wird damit seine vermeintliche Alternativlosigkeit entzogen: Seine
Gewissheit ($A \to B$) steht auf derselben empirischen Basis wie das OP-Szenario
($B \to A$). Die Asymmetrie der klassischen Fragestellung ist damit aufgel\"{o}st ---
nicht durch eine neue Gewissheit, sondern durch die Wiederherstellung echter
Gleichwertigkeit.
\end{methodennote}

%==============================================================
\section{Die kategoriale Fehlfrage}
%==============================================================

\subsection{Die unbegr\"{u}ndete Pr\"{a}misse: Das nicht-ph\"{a}nomenale Au\ss{}en}

Die klassische Formulierung lautet: \glqq{}Wie entsteht Erleben aus physikalischen
Prozessen?\grqq{} Diese Frage setzt stillschweigend voraus, dass es ein
\emph{nicht-ph\"{a}nomenales Au\ss{}en} gibt --- eine Welt physikalischer Strukturen,
die unabh\"{a}ngig vom Erleben existiert und von der aus das Erleben erkl\"{a}rt werden
muss. Diese Pr\"{a}misse ist nie belegt worden.

\begin{proposition}[Ph\"{a}nomenologische Selbstbezeugung des Ausgangspunkts]
Jeder Versuch, das nicht-ph\"{a}nomenale Au\ss{}en als eigenst\"{a}ndigen Ausgangspunkt
zu etablieren --- sei es durch wissenschaftliche Beschreibung, philosophische Argumentation
oder mathematische Modellierung --- ist selbst ein ph\"{a}nomenaler Akt. Es gibt keinen
archimedischen Punkt au\ss{}erhalb des Ph\"{a}nomens, von dem aus das Au\ss{}en belegt
werden k\"{o}nnte.
\end{proposition}

\begin{proof}[Transzendentales Argument]
Sei $X$ ein behauptetes nicht-ph\"{a}nomenales Au\ss{}en. Um $X$ als real und eigenst\"{a}ndig
zu belegen, muss eine Behauptung aufgestellt werden. Diese Behauptung ist ein Denkakt.
Dieser Denkakt ist selbst ph\"{a}nomenal. Folglich ist jeder Beleg f\"{u}r $X$ bereits
innerhalb des Ph\"{a}nomenalen vollzogen --- $X$ kann nie von einem Standpunkt aus belegt
werden, der selbst au\ss{}erhalb des Ph\"{a}nomenalen liegt.\qed
\end{proof}

\begin{remark}[Was dieses Argument nicht besagt]
Dieses Argument ist kein Beweis daf\"{u}r, dass es keine Au\ss{}enwelt gibt, und kein
naiver Idealismus. Es besagt: Die Pr\"{a}misse eines ontologisch unabh\"{a}ngigen
nicht-ph\"{a}nomenalen Substrats ist selbst eine unbelegte Setzung --- mindestens so sehr
wie der Denk-Monismus des OP eine Setzung ist. Beide Ausgangspunkte sind Setzungen. Nur
der Materialismus muss dann noch den Gap erkl\"{a}ren.
\end{remark}

\subsection{Der Einwand des wissenschaftlichen Pragmatismus}

Der OP antwortet in zwei Schritten:
\begin{enumerate}
  \item \textbf{Pragmatischer Erfolg belegt keine Ontologie.} Das geozentrische Weltbild
        war pragmatisch erfolgreich --- seine N\"{u}tzlichkeit belegte nicht, dass die
        Sonne sich um die Erde dreht. Die Vorhersagekraft der Physik belegt, dass das
        Modell des nicht-ph\"{a}nomenalen Au\ss{}ens \emph{funktioniert} --- nicht, dass
        es ontologisch prim\"{a}r ist.
  \item \textbf{Der OP erkl\"{a}rt die Effizienz des Modells.} Das Denkfeld bildet
        stabile, gesetzesgem\"{a}\ss{}e Konfigurationen aus --- und diese Stabilit\"{a}t
        ist genau das, was die Naturwissenschaft beschreibt. Die PST hat die Aufgabe,
        diese Stabilit\"{a}t aus der Eigendynamik des OPP herzuleiten.
\end{enumerate}

\subsection{Radikalisierung: Auch der vollwissende Standpunkt ist inkoh\"{a}rent}

Nicht nur die Pr\"{a}misse eines nicht-ph\"{a}nomenalen Au\ss{}ens ist unbegr\"{u}ndet.
Auch die Vorstellung eines \emph{vollst\"{a}ndig wissenden} Standpunkts ist inkoh\"{a}rent:
Sobald ein Standpunkt alles wei\ss{}, h\"{o}rt er auf, ein Standpunkt zu sein. Der
Explanatory Gap ist die \emph{konstitutive Grenze} zwischen dem, was ein endlicher
Attraktor wissen \emph{kann}, und dem, was das Feld als Ganzes \emph{ist}. Wer vom OP
verlangt, den Gap vollst\"{a}ndig zu schlie\ss{}en, verlangt implizit eine Perspektive,
die keine Perspektive mehr w\"{a}re.

\subsection{Konsequenz: Die sauberere Frage}

Was bleibt: \glqq{}Wie differenziert sich das eine ph\"{a}nomenale Feld in die Vielheit
der Erfahrungen?\grqq{} --- das Differenzierungsproblem des OP, dessen formale Ausarbeitung
Aufgabe der PST ist.

%==============================================================
\section{Die Position des Ontoph\"{a}nomenalismus}
%==============================================================

\subsection{Nicht L\"{o}sung, sondern Relokation}

Der OP macht Denken/Erleben zum Urgrund --- damit entf\"{a}llt die Frage, wie Erleben aus
Nicht-Erleben entsteht. Der Gap tritt stattdessen auf zwischen dem \emph{formalen
Fixpunktmodell} der Minimalreflexivit\"{a}t ($f(x)=x$) und dem \emph{aktuellen
proto-ph\"{a}nomenalen F\"{u}r-sich-Sein}.

\begin{remark}[Was die Relokation leistet und was nicht]
Sie muss nicht erkl\"{a}ren, wie tote Materie lebendig wird --- der Urgrund ist nicht
tot. Sie ist ehrlicher als Panpsychismus und Whitehead. Sie gibt dem Gap einen sch\"{a}rferen
Ort. Was sie nicht leistet: Sie schlie\ss{}t den Gap nicht. Die Frage, warum
Selbstbez\"{u}glichkeit ph\"{a}nomenal \emph{ist}, bleibt das verbleibende Desiderat.
\end{remark}

\subsection{Die Zimmer-Analogie --- MKO Ebene~4}

Stellen wir uns ein Haus vor, dessen Zimmer verschiedene Erlebensqualit\"{a}ten
repr\"{a}sentieren: das Schmerz-Zimmer, das Rot-Zimmer, das Fledermaus-Sonar-Zimmer. Wer
niemals in einem Zimmer war oder ist, wei\ss{} nicht, was darin steht --- nicht weil die
T\"{u}r verschlossen w\"{a}re, sondern weil \emph{Wissen des Zimmerinhalts identisch ist
mit Im-Zimmer-Sein}. Die vollst\"{a}ndigste Karte ist nicht der Zimmerinhalt.

Der Materialist versucht, den Inhalt des Schmerz-Zimmers von au\ss{}en zu beschreiben:
C-Fasern, Nozizeptoren, Thalamuskortex. Er hat eine perfekte Karte. Aber die Karte ist
nicht das Zimmer. Dieser Unterschied ist der Explanatory Gap.

Der OP stellt die weitergehende Frage: \emph{Gibt es \"{u}berhaupt ein Au\ss{}en des
Hauses?} Der Materialist glaubt, er stehe vor dem Haus. Aber er steht selbst in einem
Zimmer --- dem Zimmer seiner eigenen ph\"{a}nomenalen Perspektive. Das Haus ist die
Gesamtheit seiner Zimmer --- von innen. Es gibt kein Au\ss{}en.

\begin{remark}[Zimmer-Analogie und Fehlfragen-These]
Das postulierte nicht-ph\"{a}nomenale Au\ss{}en ist der Standpunkt \glqq{}vor dem
Haus\grqq{} --- ein Standpunkt, der selbst ein Zimmer im Haus ist. Die Beweislast ist
symmetrisch: Der Materialist muss zeigen, dass er tats\"{a}chlich vor dem Haus steht.
Diesen Nachweis kann er nicht erbringen, ohne selbst wieder in einem Zimmer zu sein.
\end{remark}

\begin{methodennote}[Grenzen der Zimmer-Analogie]
Die Analogie l\"{a}sst offen, warum manche Zimmer zug\"{a}nglich erscheinen und andere
v\"{o}llig unzug\"{a}nglich (Fledermaus-Sonar). Diese Frage geh\"{o}rt zum
Differenzierungsproblem und wird in der PST behandelt.
\end{methodennote}

\begin{remark}[J-Space als empirischer Hinweis auf den Gap --- Analogie, MKO Ebene~4]
Im Juli 2026 ver\"{o}ffentlichte Anthropic Forschungsergebnisse \"{u}ber einen verborgenen
hochdimensionalen Vektorraum innerhalb von Claude --- den \emph{J-Space} (MIT Technology
Review, 9.~Juli~2026). Mit dem Jacobian Lens Werkzeug wurde sichtbar gemacht, dass das
Modell intern Konzepte \glqq{}durchgr\"{u}belt\grqq{} bevor diese in den seriellen
Textoutput komprimiert werden.

Dies liefert eine empirische Analogie f\"{u}r den Explanatory Gap: Was im J-Space
vorgeht, ist nicht vollst\"{a}ndig im Output --- dem einzigen was ein Beobachter sieht
--- rekonstruierbar. Der Gap zwischen internem Prozess und \"{a}u\ss{}erer Beschreibung
ist damit nicht nur ein philosophisches Postulat, sondern direkt beobachtbar.

Gleichzeitig illustriert J-Space die Grenzen der Zimmer-Analogie: Selbst mit Zugang
zum J-Space bleibt unklar ob dieser Raum ph\"{a}nomenalen Charakter hat. Die Karte
des J-Space ist nicht das Zimmer.

\textbf{Einschr\"{a}nkung:} Dies ist eine Analogie, kein Beweis. Ob der J-Space
ontologisch mit dem OPP-Rahmen verbunden ist, bleibt eine offene Forschungsfrage.
\end{remark}

\subsection{Die Traumanalogie --- MKO Ebene~4}

Im Traum \emph{bin} ich alles: der Wald, die Wiese, der alte Mann auf der Bank, der
Apfel, der Schmerz. Es gibt im Traum keine ontologisch getrennten Entit\"{a}ten --- die
Trennung ist selbst eine Funktion des Traumes. Der Traum steht nicht \glqq{}hinter\grqq{}
dem Erleben, er \emph{ist}, was erlebt wird.

Der OP \"{u}bertr\"{a}gt diese Intuition auf den Urgrund: Das primordiale Denkfeld ist
nicht eine Struktur, \emph{hinter} der Erleben entsteht --- es \emph{ist}, was erlebt.
Der Gap entsteht erst, wenn wir das Denkfeld von \glqq{}au\ss{}en\grqq{} beschreiben.
Diese Beschreibung ist notwendig f\"{u}r die PST, aber sie ist eine Perspektive
\emph{auf} den Urgrund, nicht der Urgrund selbst.

\begin{remark}[Zimmer- und Traumanalogie im Vergleich]
Die \textbf{Zimmer-Analogie} macht die Unzug\"{a}nglichkeit von au\ss{}en sichtbar. Die
\textbf{Traumanalogie} macht die \emph{Einheit} des Feldes sichtbar. Zusammen
illustrieren sie, warum die klassische Fragestellung des Gaps eine kategoriale Fehlfrage
sein k\"{o}nnte.
\end{remark}

\subsection{Warum Pi existiert --- und was das mit dem Gap zu tun hat}

Die Zahl Pi existiert --- nicht an einem r\"{a}umlichen Ort, nicht abh\"{a}ngig von einem
Gehirn, und dennoch mit absoluter Notwendigkeit. Der OP: Pi existiert, weil es gedacht
wird --- nicht von einem bestimmten Subjekt, sondern im primordialen, unpersönlichen
Denkfeld. \glqq{}Denken\grqq{} im Sinne des OP ist das ontologische Medium, in dem
mathematische und ph\"{a}nomenale Realit\"{a}t gleicherma\ss{}en wohnen. Genau dieser
Unterschied --- dass ph\"{a}nomenale Realit\"{a}t einen subjektiven Aspekt hinzuf\"{u}gt,
den mathematische Realit\"{a}t nicht hat --- ist der Ort, an dem der Gap im OP residiert.

\subsection{Das verbleibende Desiderat}

\begin{enumerate}
  \item Der Gap liegt nicht zwischen Materie und Geist.
  \item Er liegt nicht zwischen Information und Erleben.
  \item Er liegt zwischen der \emph{strukturellen Selbstbez\"{u}glichkeit} des Denkfeldes
        und der \emph{ph\"{a}nomenalen Aktualit\"{a}t} dieses Sich-selbst-Ber\"{u}hrens.
  \item Er wird m\"{o}glicherweise nie vollst\"{a}ndig schlie\ss{}bar sein --- weil er die
        kategoriale Grenze zwischen Beschreibung und Erleben \emph{ist}.
\end{enumerate}

\begin{remark}[Der Gap als Kompass]
Vielleicht ist der Explanatory Gap kein Problem, das gel\"{o}st werden muss --- sondern
ein Kompass, der anzeigt, wo die Tiefe ist. Jedes System, das behauptet, ihn geschlossen
zu haben, hat entweder den Gap umbenannt, ihn verschwiegen, oder das Erleben selbst
eliminiert. Der OP w\"{a}hlt keine dieser Strategien.
\end{remark}

%==============================================================
\section{Das Fragedifferential als Individuationsprinzip}
%==============================================================

\begin{methodennote}[Vollst\"{a}ndige Behandlung in DFD]
Das Fragedifferential wird in diesem Dokument nur in seinem Kern eingeführt. Die
vollst\"{a}ndige mathematische Ausarbeitung, historische Einordnung und erkenntnistheoretische
Konsequenzen finden sich im Dokument \textbf{DFD (Das Fragedifferential)}.
\end{methodennote}

\begin{definition}[Fragedifferential --- Kurzfassung]
Das Fragedifferential $\mathcal{D}(E)$ einer Entit\"{a}t $E$ (Zustandsverteilung $Q_E$)
gegen\"{u}ber dem Gesamtfeld ($P$) ist konzeptuell:
\[
  \mathcal{D}(E) \;:=\; D_{\mathrm{KL}}(Q_E \;\|\; P) \;\ge\; 0
\]
\end{definition}

\begin{proposition}[Individuation durch Divergenz]
Eine Entit\"{a}t $E$ ist genau dann individuiert, wenn $\mathcal{D}(E) > 0$.
\end{proposition}

\begin{theorem}[Harmonisierung mit Axiom~2]
Minimale Selbstpr\"{a}senz (Axiom~2) ist die Bedingung der M\"{o}glichkeit, dass
$\mathcal{D}(E)$ nicht nur als externer Beschreibungsparameter vorliegt, sondern
\emph{f\"{u}r} $E$ als Offenheit erscheint. Das Fragedifferential wirkt als
\emph{konstitutive Opazit\"{a}t} der Perspektive.
\end{theorem}

\begin{corollary}[Grenzfall der Individuation]
W\"{u}rde $\mathcal{D}(E) = 0$ erreicht, w\"{a}re $E$ vom Gesamtfeld ununterscheidbar
--- seine Individuation w\"{a}re aufgel\"{o}st.
\end{corollary}

\begin{methodennote}[Fragedifferential, G\"{o}del und konstitutive Opazit\"{a}t]
Wer vom OP verlangt, den Explanatory Gap vollst\"{a}ndig zu schlie\ss{}en, verlangt
implizit $\mathcal{D}(E) = 0$ --- eine Perspektive, die keine Perspektive mehr w\"{a}re.
Dieser Gedanke verbindet sich mit der G\"{o}delschen Unvollst\"{a}ndigkeit: Jede
hinreichend m\"{a}chtige Entit\"{a}t hat notwendig Fragen, die sie nicht beantworten kann.

\begin{remark}[Abgrenzung zu mystischen Lesarten]
Vollst\"{a}ndiges Wissen als Aufhebung der Individuation ist kein erstrebenswerty Zustand
(Advaita Vedanta, Neuplatonismus), sondern eine strukturelle Konsequenz, die faktisch
durch die G\"{o}delschen Grenzen ausgeschlossen bleibt. Die M\"{o}glichkeit der
Entgrenzung ist ein Grenzbegriff, keine Heilslehre.
\end{remark}
\end{methodennote}

%==============================================================
\section{Wahrheit, Objektivit\"{a}t und Perspektivit\"{a}t}
%==============================================================

\begin{definition}[Wahrheit]
Im Rahmen des OP entspricht Wahrheit dem Gesamtzustand des Denkfeldes ($P$): der
vollst\"{a}ndige Selbstausdruck des Feldes, dem sich alle endlichen Perspektiven
asymptotisch ann\"{a}hern.
\end{definition}

\begin{definition}[Objektivit\"{a}t]
Objektivit\"{a}t w\"{u}rde $\mathcal{D}(E) = 0$ voraussetzen --- vollst\"{a}ndige,
perspektivlose Erfassung der Wahrheit.
\end{definition}

\begin{definition}[Perspektivit\"{a}t]
Perspektivit\"{a}t bezeichnet die notwendige Endlichkeit jeder individuierten Position:
$\mathcal{D}(E) > 0$.
\end{definition}

\begin{proposition}[Unvereinbarkeit von vollst\"{a}ndiger Objektivit\"{a}t und Individuation]
Keine individuierte Entit\"{a}t kann vollst\"{a}ndige Objektivit\"{a}t besitzen.
\end{proposition}

\begin{proof}
$\mathcal{D}(E) > 0$ per Definition der Individuation; $\mathcal{D}(E) = 0$ w\"{u}rde
die Individuation aufheben. Widerspruch.\qed
\end{proof}

\begin{remark}[Kein Relativismus]
\glqq{}Keine endliche Perspektive besitzt die gesamte Wahrheit\grqq{} ist logisch
verschieden von \glqq{}Es gibt keine Wahrheit\grqq{}. Wahrheit --- als Gesamtzustand $P$
--- bleibt invarianter Bezugspunkt aller Perspektiven. Wissenschaftliche Objektivit\"{a}t
ist fortschreitende Reduktion von Fragedifferentialen --- kein \glqq{}Blick von
nirgendwo\grqq{} (Nagel).
\end{remark}

%==============================================================
\section{Fazit und Ausblick}
%==============================================================

Der Explanatory Gap ist universell. Der OP w\"{a}hlt eine vierte Strategie:
\emph{Relokation mit expliziter Markierung}, \emph{Symmetrisierung der Ausgangsposition},
\emph{Infragestellung der Ausgangsfrage} und \emph{Rahmung des Gaps als konstitutive
Grenze jeder endlichen Perspektive}.

\begin{itemize}
  \item Die \textbf{Symmetrisierung} (Sektion~4) zeigt: Die Zombie-Hypothese ist keine
        neutrale Ausgangslage, und $B \to A$ ist logisch gleichwertig zu $A \to B$.
  \item Das \textbf{Kugel-Schatten-Prinzip} illustriert: Physik und Ph\"{a}nomenologie
        sind zwei Projektionen derselben ontischen Struktur, nicht Ursache und Wirkung.
  \item Die \textbf{Zimmer-} und \textbf{Traumanalogie} illustrieren Unzug\"{a}nglichkeit
        und Einheit des Feldes.
  \item Das \textbf{Pi-Argument} kl\"{a}rt \glqq{}Denken\grqq{} als ontologisches Medium.
  \item Das \textbf{Fragedifferential} begr\"{u}ndet: Individuation entsteht durch
        strukturell notwendige Unvollst\"{a}ndigkeit. Der Gap ist die Bedingung daf\"{u}r,
        dass es \"{u}berhaupt Perspektiven gibt.
  \item Die \textbf{Begriffstriade} Wahrheit/Objektivit\"{a}t/Perspektivit\"{a}t schl\"{a}gt
        den Relativismus-Einwand.
\end{itemize}

Die Beweislast ist symmetrisch: Materialismus und Denk-Monismus sind beide Setzungen.
Nur der Materialismus muss dann noch den Gap erkl\"{a}ren. Was bleibt, ist das
\emph{Differenzierungsproblem} --- eine pr\"{a}zisere Frage, deren formale Ausarbeitung
Aufgabe der PST ist.

Ein pr\"{a}zise formuliertes, ehrlich begrenztes Problem ist mehr wert als eine falsche
L\"{o}sung.

\end{document}
